\scanpage{54}
% L01
tel que \deleted{$g\circ\varphi_S$} l'on ait
% L02-04
\[
\left\{\begin{aligned}
g\circ\varphi_S&=\varphi_S\circ\underline{\unsure{\widetilde{g}|\widehat{g}}}\\
g\circ\varphi_A&=\varphi_A\circ\underline{\widetilde{g}}\\
g\circ\varphi_F&=\varphi_F\circ\underline{\widetilde{g}}
\end{aligned}\right.
\]

% L05-15
\uline{NB} l'une quelc. de ces trois relations implique les deux autres. Dans le cas affine, $\widetilde{g}$ défini de façon unique par $g$, et ainsi $\unsure{G|C}$ opère sur $E$. C'est encore vrai dans le cas projectif, \uline{pourvu que} $\exists\underline{s}\in\underline{S}$ avec $\varphi_S(\underline{s})\deleted{\unsure{\varphi|\ne|s}}=s$ non coplanaire avec trois des quatre pts $s_0\ s_1\ s_2\ s_3$ (unique exception: \deleted{tétraèdre à \unsure{imag.} \unsure{cop}} $S=\{s_0,s_1,s_2,s_3\}$ \ldots)
% slanted
\margin{à vérifier}
Si on veut éliminer ce \unsure{canular|cas-là}, il faut se \uline{donner} $\underline{g}\mapsto g$ faisant opérer $G$ sur $E$ \ldots
% Both ellipses above are visible manuscript continuation dots.

% L16-18
Dorénavant on suppose $\underline{X}$ régulier, on choisit repère $\underline{r}_0=(\underline{s}_0,\underline{a}_0,\underline{f}_0)$, ainsi que $\underline{\sigma}_0,\underline{\sigma}_1,\underline{\sigma}_2$ de $\underline{G}=\mathrm{Aut}(\underline{X})$, définis par
\deleted{$\unsure{\underline{\sigma}_0|\underline{r}_0}$}
% archive sheet number: 53
